\documentclass[10pt,a4paper]{article}
\usepackage[margin=20mm,headheight=15pt,headsep=7mm,footskip=10mm]{geometry}
\usepackage{fontspec,amsmath,amssymb,mathtools,microtype}
\setmainfont{lmroman10-regular.otf}[BoldFont=lmroman10-bold.otf,ItalicFont=lmroman10-italic.otf,BoldItalicFont=lmroman10-bolditalic.otf]
\setsansfont{lmsans10-regular.otf}[BoldFont=lmsans10-bold.otf,ItalicFont=lmsans10-oblique.otf,BoldItalicFont=lmsans10-boldoblique.otf]
\setmonofont[Scale=.86]{lmmono10-regular.otf}[ItalicFont=lmmono10-italic.otf]
\usepackage{ragged2e}\usepackage{xcolor,booktabs,longtable,array,graphicx,enumitem,fancyhdr,needspace,titlesec,hyperref,xurl}
\definecolor{AtlasNavy}{HTML}{18334C}\definecolor{AtlasTeal}{HTML}{116D73}\definecolor{AtlasGray}{HTML}{52616D}
\hypersetup{colorlinks=true,linkcolor=AtlasTeal,urlcolor=AtlasTeal,pdfcreator={Complexity Atlas},unicode=true}
\urlstyle{same}\setlength{\parindent}{0pt}\setlength{\parskip}{4.8pt}\setlength{\emergencystretch}{3em}
\renewcommand{\baselinestretch}{1.08}\setcounter{secnumdepth}{0}\setlist{nosep,leftmargin=1.5em}
\titleformat{\section}{\Large\bfseries\sffamily\color{AtlasNavy}\raggedright}{}{0em}{}
\titleformat{\subsection}{\large\bfseries\sffamily\color{AtlasTeal}\raggedright}{}{0em}{}
\titleformat{\subsubsection}{\normalsize\bfseries\color{AtlasNavy}\raggedright}{}{0em}{}
\titlespacing*{\section}{0pt}{14pt}{6pt}\titlespacing*{\subsection}{0pt}{10pt}{5pt}
\pagestyle{fancy}\fancyhf{}\fancyhead[L]{\small\sffamily\color{AtlasGray}COMPLEXITY ATLAS}
\fancyhead[R]{\small\sffamily\color{AtlasGray}ENGLISH EDITION}\fancyfoot[C]{\small\thepage}\renewcommand{\headrulewidth}{.3pt}
\newcommand{\AtlasMath}[1]{\par\begingroup\centering\setbox0=\hbox{$\displaystyle #1$}\ifdim\wd0>1.4\linewidth\typeout{ATLAS WIDE FORMULA: \the\wd0}\fi\ifdim\wd0>\linewidth\resizebox{\linewidth}{!}{\box0}\else\box0\fi\par\endgroup}
\newcommand{\AtlasInlineMath}[1]{\begingroup\setbox0=\hbox{$#1$}\ifdim\wd0>.4\linewidth\endgroup\AtlasMath{#1}\else\endgroup$#1$\fi}
\pdfstringdefDisableCommands{\def\AtlasInlineMath#1{#1}}
\DeclareMathOperator{\polylog}{polylog}\DeclareMathOperator{\dist}{dist}\DeclareMathOperator{\diam}{diam}
\DeclareMathOperator{\rank}{rank}\DeclareMathOperator{\tr}{tr}\DeclareMathOperator{\diag}{diag}
\DeclareMathOperator{\prox}{prox}\DeclareMathOperator{\KL}{KL}\DeclareMathOperator{\Var}{Var}
\providecommand{\eps}{\varepsilon}\providecommand{\R}{\mathbb R}\providecommand{\E}{\mathbb E}
\providecommand{\mathbbm}[1]{\mathbf{#1}}\providecommand{\mathds}[1]{\mathbf{#1}}
\providecommand{\llbracket}{[\![}\providecommand{\rrbracket}{]\!]}
\begin{document}
\raggedright

{\LARGE\bfseries\sffamily\color{AtlasNavy} Dynamic online optimization with exponentially concave losses, an unknown path budget, and single-gradient feedback.\par}\vspace{5pt}
\textbf{Setting ID.} \texttt{ONX\_\allowbreak{}DYN\_\allowbreak{}EXP\_\allowbreak{}ONEGRAD\_\allowbreak{}UNKNOWNP}

Online learning and bandits

One learner must adapt simultaneously to all comparator paths without knowing their lengths. Losses are exponentially concave, and path length is Euclidean. Each round permits only a gradient at the chosen action.

\textbf{Status.} Logarithmic gap present

\textbf{TASK tags.} formulation: Onx dyn exp onegrad unknownp; task family: Online bandit

\textbf{MODEL tags.} convexity: Exp concavity; stochastic model: Comparator path

\textbf{ORACLE tags.} interface family: Action gradient; operation: Action gradient; public information: Hidden objects explicitly declared; public information: Public inputs explicitly declared; public information: Unknown path

\textbf{ALGORITHMS tags.} adaptivity: Causal protocol; randomization: Private randomness allowed; restriction: No additional span restriction explicitly imposed; restriction: No additional zero-respecting restriction explicitly imposed

\textbf{OUTPUT tags.} certificate: No additional output certificate required; object: Online trajectory

\textbf{METRIC tags.} criterion: Dynamic regret

\textbf{GUARANTEE tags.} randomness: Algorithm private seed; randomness: Oracle samples

\textbf{RESOURCE tags.} accounting: Joint coordinates retained; measure: Native feedback; measure: Rounds; setup: Initialization included

\needspace{7\baselineskip}\section{1. Setting specification}
\needspace{4\baselineskip}\subsection{Task}
One learner must adapt simultaneously to all comparator paths without knowing their lengths. Losses are exponentially concave, and path length is Euclidean. Each round permits only a gradient at the chosen action.

\needspace{4\baselineskip}\subsection{Admissible model}
The public compact convex set \AtlasInlineMath{W\subset\mathbb{R}^d} has Euclidean diameter at most \AtlasInlineMath{D}. Losses are \AtlasInlineMath{\eta}-exp-concave on \AtlasInlineMath{W} and satisfy \AtlasInlineMath{\|\nabla f_t(w)\|_2\le G}. Arbitrary admissible sequences are allowed.


\AtlasMath{e^{-\eta f_t(\cdot)}\ \text{\text{on} }W\text{ \text{is concave}}.}


Precisely, exponential concavity requires that in each round \AtlasInlineMath{t}, the map \AtlasInlineMath{x\mapsto\exp(-\eta f_t(x))} be concave on the stated feasible domain. This property concerns the complete action variable; the norm and gradient bound are specified by this axis.

\needspace{4\baselineskip}\subsection{Oracle interface}
In each round, the learner first commits to its final action \AtlasInlineMath{w_t\in W} and then receives only \AtlasInlineMath{\nabla f_t(w_t)}. Full-function feedback and additional value or gradient evaluations at other points are unavailable.

Exact Mahalanobis projection onto the public set is allowed. The arithmetic cost of one projection is \AtlasInlineMath{C_{\mathrm{proj}}(d)}.

The learner is causal and remains in the domain. Public inputs are \AtlasInlineMath{d,T,D,G,\eta,W}; comparator path budget \AtlasInlineMath{P_T^{(2)}} is not supplied.

Received feedback may be stored and reused in computation. Unrevealed information cannot be requested from past rounds, and consumed random samples cannot receive additional queries. Current-round exceptions are limited to the explicitly supplied full-function access or same-sample two-point interface.

\needspace{4\baselineskip}\subsection{Output}
The output is the complete action trajectory satisfying feasibility and timing requirements. Randomized mixtures may be implemented before acting. Retrospective averages or convex combinations cannot replace actions that have already incurred loss. No additional optimal value or certificate is required.

\needspace{4\baselineskip}\subsection{Loss or accuracy criterion}

\AtlasMath{P_T^{(2)}(u)=\sum_{t=2}^{T}\|u_t-u_{t-1}\|_2,\qquad \operatorname{Reg}_T(u)=\sum_{t=1}^{T}\bigl(f_t(w_t)-f_t(u_t)\bigr).}



\AtlasMath{\Lambda_T=1+\log(2T),\qquad P_T^{(2)}\le \frac{cT}{\sqrt d\,\Lambda_T^2}.}


\needspace{4\baselineskip}\subsection{Resource accounting}

\AtlasMath{N=T,\qquad\text{Total arithmetic cost}=O\!\left(T(d^3+C_{\rm proj}(d))\log^2T\right).}



\AtlasMath{\text{Storage}=O(d^2\log^2T).}


Storage counts real registers. Gradient-call count excludes internal matrix computations and projections.

Exploration, initialization, and system identification all consume the stated online budget. Queries, arithmetic, communication rounds or bits, and memory remain joint resources. No arithmetic or bit-complexity optimality is asserted.

\needspace{4\baselineskip}\subsection{Algorithm views and success guarantees}
All measurable causal randomized algorithms obeying the stated information protocol are allowed. No linear-span, zero-support, or specific iteration restriction is added implicitly. Internal computation and memory are unrestricted except for the explicit computation, communication, or privacy constraints.

Private randomization is allowed, with the seed independent of the fixed instance and future external samples. Internal arithmetic, memory, linear spans, and iteration forms are unrestricted except for the stated computation, communication, privacy, and within-round interaction limits.

One algorithm gives the pathwise exponentially concave branch bound for comparator paths within its parameter window. The minimum frontier across branches under full-function feedback cannot be transferred directly to this interface.

Here \AtlasInlineMath{c>0} is the constant used by the corresponding atomic upper bound. In dimension-dependent comparisons, \AtlasInlineMath{D,G,\eta} is fixed and does not vary with \AtlasInlineMath{d,T}.

\needspace{7\baselineskip}\section{2. Notation and parameters}
\AtlasInlineMath{D}: A Euclidean-diameter bound on the action set. \textbf{Domain:} \AtlasInlineMath{D\ge0}.

The condition is \AtlasInlineMath{\sup_{x,y\in X}\|x-y\|_2\le D}; replace X by W when W denotes the action set.

\AtlasInlineMath{G}: Uniform Euclidean bound on the full loss gradient in each round \textbf{Domain:} \AtlasInlineMath{G>0}.

The condition \AtlasInlineMath{\|\nabla f_t(x)\|_2\le G} holds for all t and x specified by the model.

\AtlasInlineMath{T}: Online rounds within the fixed evaluation horizon. \textbf{Domain:} \AtlasInlineMath{T\in\mathbb N,\ T\ge1}.

The information protocol fixes feedback count per round. The horizon \AtlasInlineMath{T} is not an independent cumulative scalar-query count.

\AtlasInlineMath{d}: The ambient dimension of actions or feature vectors. \textbf{Domain:} \AtlasInlineMath{d\in\mathbb N,\ d\ge1}.

Vectors belong to \AtlasInlineMath{\mathbb R^d}; dimension is still denoted d when a non-Euclidean norm is used.

\AtlasInlineMath{\eta}: Exponential-concavity coefficient of the loss \textbf{Domain:} \AtlasInlineMath{\eta>0}.

The function \AtlasInlineMath{x\mapsto\exp(-\eta f_t(x))} is concave on the action set. This parameter is not an algorithmic step size.

\AtlasInlineMath{P}: Euclidean comparator-path budget used only for stratified evaluation \textbf{Domain:} \AtlasInlineMath{0\le P\le D(T-1)}.

The comparator satisfies \AtlasInlineMath{P_T^{(2)}(u)=\sum_{t=2}^T\|u_t-u_{t-1}\|_2\le P}. A common algorithm may not tune to \AtlasInlineMath{P}.

\AtlasInlineMath{W}: Public nonempty compact convex action set \textbf{Domain:} \AtlasInlineMath{W\subset\mathbb R^d}.

\AtlasInlineMath{\operatorname{diam}_2(W)\le D}.

\AtlasInlineMath{w_t,u_t}: Learner action and current comparator location \textbf{Domain:} \AtlasInlineMath{w_t,u_t\in W}.

\AtlasInlineMath{f_t}: Current differentiable loss \textbf{Domain:} A real-valued function on W.

The function \AtlasInlineMath{\exp(-\eta f_t(\cdot))} is concave, with Euclidean gradient norm at most G.

\AtlasInlineMath{P_T^{(2)}(u)}: Realized total Euclidean path length of the comparator \textbf{Domain:} \AtlasInlineMath{0\le P_T^{(2)}\le D(T-1)}.

\AtlasInlineMath{P_T^{(2)}(u)=\sum_{t=2}^T\|u_t-u_{t-1}\|_2}.

\AtlasInlineMath{\operatorname{Reg}_T(u)}: Dynamic regret relative to the complete comparator path \textbf{Domain:} Real number.

\AtlasInlineMath{\operatorname{Reg}_T(u)=\sum_t[f_t(w_t)-f_t(u_t)]}. One algorithm satisfies the upper bound simultaneously for all comparator paths.

\AtlasInlineMath{\Lambda_T}: Logarithmic factor appearing in the source upper bound \textbf{Domain:} Positive real number.

\AtlasInlineMath{\Lambda_T=1+\log(2T)}.

\AtlasInlineMath{c_0}: A positive universal constant in the source lower-bound horizon window. \textbf{Domain:} A fixed positive constant.

Require \AtlasInlineMath{T\ge c_0d^2}; the source gives no numerical value here, so none is assigned arbitrarily.

\AtlasInlineMath{A,C_{\mathrm{proj}}(d)}: The positive-definite metric for Mahalanobis projection and the cost of one exact projection. \textbf{Domain:} Positive definite matrix and arithmetic-operation count.

Metric projection minimizes \AtlasInlineMath{(x-y)^{\mathsf T}A(x-y)}; computational cost and loss-feedback count are measured separately.

\AtlasInlineMath{X}: Alias for the action set in the lower-bound parameter slice \textbf{Domain:} An alias for W.

\AtlasInlineMath{X=W=B_2^d(1/2)}.

\AtlasInlineMath{c}: A fixed positive constant in the exponentially concave comparison window. \textbf{Domain:} Positive real number.

Applies when \AtlasInlineMath{P\le cT/(\sqrt d\,\Lambda_T^2)}. Preserve the source's constant range without assigning numerical values.

\AtlasInlineMath{\mathcal A}: Learning policies satisfying the specified information, computation, and privacy restrictions \textbf{Domain:} A measurable causal policy.

Actions may depend only on revealed information and the private seed. Additional restrictions are governed by the algorithm protocol in this profile.

\AtlasInlineMath{t}: Interaction-time index \textbf{Domain:} \AtlasInlineMath{t\in\{1,\ldots,T\}}.

\AtlasInlineMath{\mathbb E}: Mean operator \textbf{Domain:} Expectation over the specified random variables.

Expectation includes the learner's private randomness and, when allowed, the environment strategy's randomness. A fixed loss sequence introduces no additional observation noise.

\AtlasInlineMath{\log,\log_2}: Natural and binary logarithms \textbf{Domain:} Function on the positive real axis.

\AtlasInlineMath{\log x=\ln x}, \AtlasInlineMath{\log_2x=\log x/\log2}.

\AtlasInlineMath{c,C}: Positive constant in the bound without an assigned numerical value \textbf{Domain:} Positive real number.

Constants apply only under the associated bound’s conditions and fixed-parameter convention. Unspecified theorem constants are not treated as additional algorithm inputs.

\AtlasInlineMath{O,\Omega,\widetilde O,\widetilde\Omega}: Asymptotic upper and lower bounds, including explicitly permitted logarithmic suppression. \textbf{Domain:} Asymptotic comparison notation.

The result applies only in the regime declared for that bound. Tilde notation hides only its stated logarithms and fixed-parameter dependencies, never additional polynomial parameter factors.

\AtlasInlineMath{N}: Cumulative observations counted according to the declared feedback type. \textbf{Domain:} Nonnegative integer or integer-valued stopping time.

The quantity \AtlasInlineMath{N=T} counts full-gradient or subgradient responses. Projection, proximal, and matrix operations are counted separately.

\needspace{12\baselineskip}\section{3. Complexity bracket}
\needspace{4\baselineskip}\subsection{Compared quantities}
\needspace{4\baselineskip}\subsubsection{Q1 - native performance}
\AtlasMath{P_T^{(2)}(u)=\sum_{t=2}^{T}\|u_t-u_{t-1}\|_2,\qquad \operatorname{Reg}_T(u)=\sum_{t=1}^{T}\bigl(f_t(w_t)-f_t(u_t)\bigr).}
\textbf{Scope.} Compare only within this setting and the explicit regime, retaining multiple objectives and time normalization.

\textbf{Optimization direction.} minimize

\needspace{4\baselineskip}\subsection{R1 - SUPPORTED}
\textbf{Region.} Fix \AtlasInlineMath{X=B_2^d(1/2)} and \AtlasInlineMath{D=1,G\ge 3/2,0<\eta\le 4/9}. Assume \AtlasInlineMath{T\ge c_0d^{2}} and \AtlasInlineMath{P\le \min{T/d^{2},cT/(\sqrt{d}\Lambda_T^{2})}}.

\textbf{Quantity.} Q1

\textbf{Resource.} Fixed native online budget, with computational, communication, and privacy resources retained separately.

\textbf{Status.} Logarithmic gap

\textbf{Lower bound.} B2

\AtlasMath{\Omega(d+d^{1/3}T^{1/3}P^{2/3})}
\textbf{Upper bound.} B1

\AtlasMath{\widetilde{O}(d+d^{1/3}T^{1/3}P^{2/3})}
\needspace{7\baselineskip}\section{4. Bounds and supporting evidence}
The following conditions apply to every bound in this profile, in addition to each bound's stated region and algorithm scope.

\textbf{Oracle mapping.} Retain the specified feedback type, sample coupling, public information, and environment permissions.

\needspace{4\baselineskip}\subsection{B1 - A deterministic learner adapting to curvature using one gradient.}
\textbf{Admission.} Included in the bracket

\textbf{Direction.} upper

\textbf{Witness or lower-bound scope.} Deterministic witness

\AtlasMath{O\!\left(d\Lambda_T^2+d^{1/3}\Lambda_T^{4/3}T^{1/3}\bigl(P_T^{(2)}\bigr)^{2/3}\right)}
\textbf{Sources.} \href{https://arxiv.org/pdf/2610.20187v1}{[1]}

\textbf{Conditions.} Assume \AtlasInlineMath{\Lambda_T=1+\log(2T)} and \AtlasInlineMath{P_T^{(2)}\le cT/(\sqrt d\,\Lambda_T^2)}. The quantity \AtlasInlineMath{c>0} is an absolute constant left unspecified in the source.; The quantity \AtlasInlineMath{D,G,\eta} is fixed independently of dimension. This expression gives only the single-gradient branch in its window, not the full minimum over several branches.

\textbf{Source location.} \href{https://arxiv.org/pdf/2610.20187v1}{[1]} Theorem 4.1 p.6 Eq.(5)

\needspace{4\baselineskip}\subsection{B2 - Block construction with squared losses; lower bound for arbitrary randomized learners}
\textbf{Admission.} Included in the bracket

\textbf{Direction.} lower

\textbf{Witness or lower-bound scope.} All randomized adaptive algorithms

\AtlasMath{\Omega\!\left(d+\min\!\left\{d^{1/3}T^{1/3}P^{2/3},\ \sqrt{TP},\ T\right\}\right)}
\textbf{Sources.} \href{https://arxiv.org/pdf/2610.20187v1}{[1]}

\textbf{Conditions.} The quantity \AtlasInlineMath{P\ge0} is an analytical upper bound on the comparator's total Euclidean path length and need not be supplied to the learner. Assume \AtlasInlineMath{T\ge c_0d^2}, with dimension-independent \AtlasInlineMath{c_0>0}.; Use action set \AtlasInlineMath{W=B_2^d(1/2)} and normalizations \AtlasInlineMath{D=1}, \AtlasInlineMath{G\le3/2}, and \AtlasInlineMath{\eta=4/9}, which also satisfy the weaker requirement \AtlasInlineMath{0<\eta\le4/9}. Construct losses \AtlasInlineMath{f_{j,y}(w)=(w_j-y)^2/2}, where \AtlasInlineMath{1\le j\le d} and \AtlasInlineMath{y\in\{-1,1\}}.; This is a worst-case lower bound on the specified normalized domain. It does not apply uniformly to every fixed compact convex set.

\textbf{Derivation.} A lower bound covering all randomized algorithms with full-function feedback also applies to the weaker interface returning only the gradient at the current action. Preserve the radius-\AtlasInlineMath{1/2} ball, the source curvature normalization, and condition \AtlasInlineMath{T\ge c_0d^2}.

\textbf{Source location.} \href{https://arxiv.org/pdf/2610.20187v1}{[1]} Theorem 5.1, PDF pp.8-9, Eq.(8)

\needspace{7\baselineskip}\section{References}
[1] Yanyi Li; Wendao Wu; Haihan Zhang; Qinshi Wang; Chenheng Zhang; Chunyuan Zheng; Cong Fang; Haoxuan Li; Zhouchen Lin. Dimension-Optimal Dynamic Regret in Exp-Concave Online Learning. 2026. 2610.20187v1.

\url{https://arxiv.org/pdf/2610.20187v1}

\end{document}
